% ============================================
% 第11章：教师是"改稿机器"吗？——AI时代的教育伦理
% 悬索桥 · 教育伦理
% ============================================

\chapter{教师是``改稿机器''吗？}

\vspace{0.5cm}

\begin{center}
{\small\color{muted} 悬索桥的主缆，承受着整座桥的重量。}
{\small\color{muted} 但如果锚碇松了，主缆再强也没有用。}
{\small\color{muted} 教师的角色，就是教育的锚碇。}
\end{center}

\vspace{1cm}

\section{当AI可以批改一切}

2026年6月23日，我在跟两位老师讨论AI教学平台。我讲到了平台上的AI批改功能：

\begin{shiyu-inline}
``AI会去过滤，说你这个评论是不是在灌水。有的学生最开始第一版的时候，他打一一一一，凑够十个一，我这要求至少十个字嘛，一点提交，叭就过去了，他的分拿到手了。但是这种情况非常少，更多的学生会认真的写。''

\hfill ——2026年6月23日，AI教学平台讨论，录音转写
\end{shiyu-inline}

AI可以自动检测灌水评论。AI可以自动聚合三十个人的评语。AI可以自动判断一个学生的参与是否``有效''。AI可以自动给分。

这些功能，在传统的课堂里，都是老师做的。老师要花几个小时看评论、找灌水、打分、汇总。现在，AI几秒钟就做完了。

\textbf{然后呢？}

老师多出来的时间，用来做什么？用来设计更好的困惑？用来跟学生一对一的交流？用来反思自己的教学？还是——用来变成一台``改稿机器''？

\section{``改稿机器''}

2026年6月12日，师宇在复盘一场铁路局AI培训时，提到了一个让我警觉的场景：

\begin{shiyu-inline}
``培训现场也触及深层焦虑：AI可能让人不会思考，导师、老师、管理者可能被迫变成'改稿机器'。''

\hfill ——师宇博文，2026年6月12日
\end{shiyu-inline}

\textbf{改稿机器。}

学生用AI生成初稿。老师用AI检测AI。老师改几处。学生再生成。老师再检查。循环往复。老师和学生之间，不再有真实的交流。只有AI生成的内容，在两个屏幕之间来回传输。

这可能是AI时代教育最可怕的未来：\textbf{不是AI取代了教师。是教师变成了AI的质检员。}

\section{什么必须保留给人？}

2026年7月，在教学理念综合分析中，我列出了\textbf{不可外包给AI的环节}：

\begin{shiyu-inline}
``保留'不可外包给AI'的环节：自我判断懂没懂、解释关键选择、失败复盘和迁移总结由学生完成。''

\hfill ——教学理念综合分析，2026年7月
\end{shiyu-inline}

这不是一个技术问题。这是一个伦理问题。

\textbf{自我判断懂没懂。}AI可以给你一个自测题，告诉你得分。但它不能替你感受``我懂了''那个瞬间。那个瞬间，必须是你自己经历的。

\textbf{解释关键选择。}AI可以解释它为什么选这个方案。但那个解释是统计的——``因为数据里这个方案最常见''。那不是解释。那是概率。真正的解释需要你理解：在什么约束下，我为什么放弃了A选了B。

\textbf{失败复盘。}AI可以分析你的错误模式。但它不能替你感受失败——那种``我明明知道但就是做错了''的懊悔，那种``我下次一定不会再犯''的决心。这些感受，是学习最核心的驱动力。

\textbf{迁移总结。}AI可以帮你总结``这次学到了什么''。但它总结的是它看到的模式。你真正学到的东西，可能和AI总结的完全不同。因为\textbf{学习是转化}——你被改变了，你变成了另一个人。AI看不到这种改变。

\section{教师的新角色}

如果AI能做知识传递，能做作业批改，能生成教案和PPT——那教师还能做什么？

\textbf{教师能做三件事，AI一件都做不了。}

\textbf{第一，口头诊断。}2026年1月9日，混凝土结构答辩。我改了一个参数，问学生：``为什么配筋量变了？''那个问题，不是在考他知不知道答案。我是在\textbf{诊断}——诊断他理解到了哪一层，诊断他的思维卡在哪里，诊断下一步该往哪个方向引导。AI可以做选择题。AI不能做口头诊断。因为口头诊断需要\textbf{看见}——看见学生的眼神，听见他的犹豫，感受他是在``回忆''还是在``推理''。

\textbf{第二，设计困惑。}AI可以生成任何内容。但它不能判断：\textbf{这个内容，在什么时间、以什么方式、给什么学生，能制造出最有效的困惑。}因为困惑是\textbf{个人的}。同一个问题，对A学生是困惑，对B学生是常识。设计困惑，需要你了解每一个学生——不是他们的分数，是他们\textbf{是谁}。

\textbf{第三，见证骄傲。}AI可以给分。但它不能\textbf{见证}。当学生指着自己搭的桥说``站住了''，那个眼神，那个语气，那个嘴角微微上扬的弧度——AI看不到。它不知道那个学生在三个月前差点转专业。它不知道那个学生曾经在深夜因为一个参数调不对而摔过键盘。它不知道那个瞬间，对那个学生意味着什么。

\textbf{见证，是教师最不能被替代的能力。}

\begin{insight}
\label{insight:chap11}

\textbf{AI时代，教师最值钱的能力，不是``讲得好''。}

是\textbf{听得懂}——听懂学生卡在哪里。是\textbf{看得见}——看见学生努力了但没做成。是\textbf{等得起}——等学生自己发现，而不是急着告诉他答案。是\textbf{记得住}——记得这个学生三个月前的样子，知道他变了多少。

这些，AI一件都做不了。
\end{insight}

\section{悬索桥的隐喻之二}

悬索桥的主缆，承受着整座桥的重量。但主缆的力，最终要传到锚碇——埋在地下，看不见，但绝对不能松。

\textbf{教师的角色，就是教育的锚碇。}

AI可以拉起主缆。AI可以生成内容、批改作业、分析数据。但所有这些力的终点，必须是一个\textbf{人}——一个能判断、能诊断、能见证、能在学生指着自己搭的桥说``站住了''的时候，真心为他高兴的人。

锚碇松了，悬索桥就塌了。教师退后了，教育就变成了两台机器之间的数据传输。

\begin{shiyu-note}
\label{note:chap11}

\textbf{不要让AI把你变成改稿机器。}

当AI帮你批改作业的时候，省下来的时间，不要用来批改更多作业。用来跟学生聊天。用来观察他们搭桥。用来问他们：``你为什么这么搭？''——然后认真听他们回答。

那些回答里，有你用任何AI都检测不出来的东西：一个正在努力理解世界的人，正在变成另一种人。
\end{shiyu-note}

\vspace{1cm}

\begin{decision-point}
\label{decision:chap11}

\noindent\textbf{这一章，我们看到了AI时代教师角色的重新定义。}

\vspace{0.3cm}

\noindent 下一节课，不要站在讲台上。走下去。走到一个学生旁边。看他搭桥。不要说话。等他自己发现错了。等他改。等他再试。等桥站住。

\noindent 然后只说一句话：``你刚才改的那一下，为什么管用？''

\noindent 听他的回答。你会听到一些AI永远听不到的东西。
\end{decision-point}

\vspace{0.5cm}

\begin{flushright}
{\small\color{muted} 下一章：桥不会自己存在——关于责任与骄傲。}
\end{flushright}